\subsection*{Управление смещением и разбросом для оптимизации моделей машинного обучения}

На самом деле, можно строить следующие параллели:
\begin{itemize}
    \item Как мы боремся с недообучением "--- ровно такие же методы мы применяем для борьбы с высоким смещением;
    \item Аналогично с высоким разбросом и переобученим.
\end{itemize}

Т.е. для <<подавления>> роста смещения можно использовать уже известные регуляризации, например. Или, как вариант, можно разбить исходную выборку на 2: на одной из выборок обучаться, а на другой "--- проверять качество модели. Если на второй выборке качество плохое "--- модель переобучена.

\subsection*{Кросс-валидация и метод KFold: обеспечение надёжности оценки моделей}

Однако у метода, который мы рассмотрели ранее, есть одна большая проблема "--- разбиение происходит случайным образом, и соответственно, мы не можем гарантировать, что данные на тестовой выборке (на которой проверяем качество модели) модель вообще будет в состоянии обработать. А что если первый, тренировочный набор содержит вообще непохожие примеры?

Соответственно, мы сильно зависим от случайности этого разбиения. Но как можно этого избежать?

Техника, которая позволяет получить более объективную оценку, называется \texttt{KFold}. Её алгоритм заключается в следующем: исходный набор данных случайным образом разбивается на $K$ равных по размеру частей, называемых фолдами (folds). Далее процесс обучения и оценки повторяется $K$ раз. На каждой итерации один из $K$ фолдов по очереди выступает в роли тестовой выборки (validation fold), а остальные $(K-1)$ фолды объединяются в обучающую выборку. Модель обучается на $K-1$ фолдах, а потом тестируется на одном фолде. Так повторяется до тех пор, пока не будут рассмотрены все такие разбиения данных на $K$ фолдов. 

Результатом работы алгоритма является массив из $K$ оценок качества. Итоговая оценка модели вычисляется как среднее арифметическое этих $K$ значений, а её стабильность можно оценить по стандартному отклонению.